\newpage

\section{Conclusiones.}
\paragraph{}
A modo de conclusión puede decirse que con con la realización de este trabajo se comprobó efectivamente la disparidad que exite entre el comportamiento de un amplificador operacional ideal y uno real. Si bien la teoría resulta útil para estimar la respuesta esperda en los dispositivo electrónicos es importante tener siempre presente que en la práctica el funcionamiento de cualquier dispositivo o componente difiere de los parámetros obtenidos mediante el análisis y cálculos teóricos. Esta discrepaciancia entre un resultado y otro puede resultar determinante al momento de elegir un amplificador operacional para la implementación de un proyecto; la exactitud y precisión deseadas para la respuesta del sistema, así como la robustez esperada son factores que deciden sobre el empleo del amplificador en cuestión en función de sus características reales.
\paragraph{}
Los datos proporcionados por el fabricante del dispositivo a utilizar permiten hacer una primera contrastación respecto a los cálculos que puedan efectuarse, sin embargo para lograr tomar una decisión certera el procedimiento ideal consite en ensayar cada uno de los posibles candidatos que se tienen para llevar a cabo la implementación. Dado que en mucho casos esto no es posible, la disponibilidad y uso de simuladores que asemejen el funcionamiento de los amplificadores de interés al elemento real resulta fundamental. Para este caso, se empleó el simulador LTSpice el cual ofrece una aproximación satisfactoria del comportamiento de componentes electrónicos permitiendo verificar que el análisis previamente realizado por los miembros de este grupo sobre el modelo de amplificador sobre el cual se realizaron los ensayos.
